% !Mode:: "TeX:UTF-8"
% !TeX program = xelatex
%%%%%%%%%%%%%%%%%%%%%%%%%%%%%%%%%%%%%%%%%%%%%%%%%%%%%%%%%%%%%%%%%%%%%%%%%%%%%%%%
%       __    _ ________              _                                        
%      / /_  (_)_  __/ /_  ___  _____(_)____          Happy LaTeXing!          
%     / __ \/ / / / / __ \/ _ \/ ___/ / ___/                                   
%    / / / / / / / / / / /  __(__  ) (__  )           Copyright (C) @dustincys 
%   /_/ /_/_/ /_/ /_/ /_/\___/____/_/____/    v3.2a   Email: yanshuoc@gmail.com
%   https://github.com/hithesis/hithesis              https://yanshuo.site     
%                                                                              
%        -: . :=:        =+ .*.       :-:*=#+:        :. --                    
%       :%#.%%*=:       .@%:%-%       :@%-=-:.       :@=*@%@:                  
%      .*@+.@=:=:     #@@%::=%:.      =%*=:%.:.   :%%@%.:=+                    
%     =@@#.-%%=-.      .+%:*= :-.    .@%+%%%%:      +@%:+: +%*-                
%      :%:+.#:*  :-   :@@. %%*@*    .#@=%#-@#.     #%%+=@%=@*                  
%     :+. *=  *==#=     *. =#%%     *%:+:#+==#*-     =: =:#*:                  
%                                        .:                :        :=:        
%   @syvshc            .:=:  .@*         %#.        =@@*.  .=       .*@--*:    
%   @MingZZX           .*@- =@%:%:    :%@@@%       .#@-@=#-::    #@*.*@*.      
%   @xiF616             #%  ** :%.     .*@#%-      -+=%  %:-:    *-.*=.*:      
%   @Xeonacid         .*%*-+%=.+=   :#%%@%#.       .#@#-.  -:     -%+%@*-.     
%   @SchrodingerBlume    .*@:  =*.     .#- :@%:     +%*.  ..**     .=*.@:.:::. 
%                                   .-:.    :==:          ..:       .*=        
%%%%%%%%%%%%%%%%%%%%%%%%%%%%%%%%%%%%%%%%%%%%%%%%%%%%%%%%%%%%%%%%%%%%%%%%%%%%%%%%
\documentclass[fontset=fandol,kind=thesis,degree-level=doctor,campus=harbin]{hithesis}
% 一份主文件管四种交付物，写哪一种由 kind 决定：
%   kind=thesis            学位论文（中文）
%   kind=thesis,lang=en    学位论文（英文）
%   kind=proposal           开题报告
%   kind=interim           中期报告
% 各交付物用哪些文件写在下面对应的 \hitsetup[kind=…] 里，不必改这一行以外的东西。
%
% 其余文档类选项的含义见手册，常用的几个：
% degree-level=doctor|master|bachelor|postdoc   必填
% degree-type=academic|professional             学术学位还是专业学位，默认 academic
% campus=harbin|shenzhen|weihai                 默认 harbin
% lang=zh|en                                    正文语言，默认 zh

\graphicspath{{figures/}}

% 示例宏包。里面是示例文档自己要用的东西：术语词条、染色体示意图的绘图宏、
% 数学记号速记、\num 的分组空格。都不是规范要求的，复制一份改成自己的。
\usepackage{hithesis}

\begin{document}

%%%%%%%%%%%%%%%%%%%%%%%%%%%%%%%%%%%%%%%%%%%%%%%%%%%%%%%%%%%%%%%%%%%%%%%%%%%%%%%%
% 共用的元信息。各自不同的部分写在对应的封面文件里。
%%%%%%%%%%%%%%%%%%%%%%%%%%%%%%%%%%%%%%%%%%%%%%%%%%%%%%%%%%%%%%%%%%%%%%%%%%%%%%%%
\hitsetup{
  info={
    student-id={9527},
    zh={
      title={局部多孔质气体静压}{轴承关键技术的研究},
      affiliation={机电工程学院},
      specialty={机械制造及其自动化},
      author={于冬梅},
      supervisor={某某某教授},
    },
  },
}

%%%%%%%%%%%%%%%%%%%%%%%%%%%%%%%%%%%%%%%%%%%%%%%%%%%%%%%%%%%%%%%%%%%%%%%%%%%%%%%%
% 各交付物的部件。方括号里是条件，只有当前这份交付物那一组会生效。
% 顺序由类按规范决定，这里只说哪个部件用哪个文件；不写的部件不排。
%%%%%%%%%%%%%%%%%%%%%%%%%%%%%%%%%%%%%%%%%%%%%%%%%%%%%%%%%%%%%%%%%%%%%%%%%%%%%%%%
% 四组条件两两互斥，任何一份文档只命中一组，所以先后无关，各组必须自足。
% （\hitsetup 是后写覆盖先写，条件只决定这一句做不做、不决定谁优先。写成互斥的
% 就不必琢磨顺序；写成「一组打底、另一组只改几项」也行，但那样顺序就有讲究了。）
\hitsetup[kind=thesis,lang=zh]{ structure = {
  frontmatter = {
    cover      = {thesis/frontmatter/cover},
    abstract   = {thesis/frontmatter/abstract},
    denotation = {thesis/frontmatter/denotation},
    contents   = true,
  },
  % 正文各章按写的顺序排，结论是正文的最后一章。
  mainmatter = {thesis/mainmatter/regu, thesis/mainmatter/introduction,
                thesis/mainmatter/recruit, thesis/mainmatter/conclusion},
  backmatter = {
    bib              = {reference},
    appendix         = {thesis/backmatter/appA},
    publications     = {thesis/backmatter/publications},
    resolution       = true,
    index            = {thesis/backmatter/ceindex},
    authorization    = true,
    acknowledgements = {thesis/backmatter/acknowledgements},
    resume           = {thesis/backmatter/resume},
  },
  bib-style = {hithesis},   % 深圳校区用 hitszthesis
} }

% 英文版：致谢在前置，没有符号表、附录、索引、决议、简历。
\hitsetup[kind=thesis,lang=en]{ structure = {
  frontmatter = {
    cover            = {thesis/frontmatter/cover-en},
    contents         = true,
    acknowledgements = {thesis/frontmatter/acknowledgements-en},
  },
  mainmatter = {thesis/mainmatter/eng_regu, thesis/mainmatter/conclusion},
  backmatter = {
    bib           = {reference-en},
    publications  = {thesis/backmatter/publications},
    authorization = true,
  },
  bib-style = {hithesis},
} }

\hitsetup[kind=proposal]{ structure = {
  frontmatter = { cover = {report/proposal/frontmatter/cover} },
  mainmatter  = {report/proposal/mainmatter/shenzhen_doctor},
} }

\hitsetup[kind=interim]{ structure = {
  frontmatter = { cover = {report/interim/frontmatter/cover} },
  mainmatter  = {report/interim/mainmatter/shenzhen_doctor},
} }

% 前置、正文、后置由类按 structure 自动排，三个 matter 宏不用写。
% 想在两段之间插点别的东西时，把 \hitfrontmatter、\hitmainmatter、
% \hitbackmatter 显式写出来即可，写了的就在写的位置排。详见手册。
\end{document}
% Local Variables:
% TeX-engine: xetex
% End:
